\subsection{The Pile 的社区意义}

The Pile 不是只是”另一个大语料”，而是开源社区把数据配方从私有资产转成公共议题的标志。它把书籍、论文、代码、百科、问答和其他域混在一起，让”开源模型的数据应该长什么样”第一次变成一个可讨论的问题。

The Pile 中几个值得关注的子集：
\begin{itemize}
    \item \textbf{Project Gutenberg}：始于 1971 年，截至 2025 年约有 7.5 万本书，主要是版权已过期的公共领域作品（莎士比亚、贝多芬等）。
    \item \textbf{Books3}：来自影子图书馆 Bibliotik 的 19.6 万本书，包含 Stephen King、Zadie Smith 等当代作者的作品。因版权侵权已被下架。
    \item \textbf{StackExchange}：用户贡献的问答数据，有声望系统和标签，Q\&A 格式天然接近指令跟随场景。
    \item \textbf{GitHub}：代码数据不仅有助于编程任务，folklore 认为也有助于推理能力。存在大量重复（fork、复制的代码）。
\end{itemize}

\begin{knowledgebox}{影子图书馆与训练数据}
Library Genesis (LibGen)、Z-Library、Anna's Archive、Sci-Hub 等影子图书馆绕过版权和付费墙，提供了海量书籍和论文。LibGen 有约 400 万本书（2019 年数据），Sci-Hub 有约 8800 万篇论文（2022 年数据）。这些资源在学术界和发展中国家有巨大需求，但在法律上争议极大。Meta 被曝出使用 LibGen 训练模型，引发了作者群体的集体诉讼。
\end{knowledgebox}

